\section{Group or zone? Where the difference comes from}\label{sec:zones}

\textbf{Question.} Group time is more ballistic than solo flights at long lags. Is this the group, or the ground where groups fly
(and where solo pilots would also progress in a straight line)? \textbf{Method.} Each piece of track is assigned to a \emph{stratum} (zone = cell
of 25~km of its middle; zone$\times$month; zone$\times$hour; site-day$\times$hour). In every stratum with both populations the ensemble MSD
of the solo flights S (and of the same flights alone A) is \emph{reweighted to the mix of the group G over the strata}: S then flies, in
the statistic, where and when G flies. We compare $H_\mathrm{loc}(\tau)$ over \textbf{every lag} and the fitted $H$. The share of G with a counterpart
in its stratum is reported; ``as is'' is the same comparison without the reweighting, on the same strata and units (a flight counts once per
cell it crosses).

\begin{figure}[H]
\centering
\includegraphics[width=\textwidth]{zones_curves.pdf}
\caption{$H_\mathrm{loc}(\tau)$ of G (orange), of S reweighted to the strata of G (blue) and of A reweighted to them (green), 5--95\% bands.
First panel: all flights pooled (the raw comparison). ``Share of G kept'': flights of G with a solo counterpart in their stratum.}
\label{fig:zones}
\end{figure}

\begin{figure}[H]
\centering
\includegraphics[width=0.92\textwidth]{zones_differences.pdf}
\caption{$\Delta H_\mathrm{loc}(\tau)$ over every lag: G minus S (left) and G minus A (right), the other population reweighted to the strata of G; one line per
stratification, finer strata in darker orange.}
\label{fig:zones-diff}
\end{figure}

\begin{table}[H]
\centering\footnotesize\setlength{\tabcolsep}{2.6pt}
\begin{tabular}{@{}l r c c c c c c c c c@{}}
\toprule
 & & \multicolumn{3}{c}{fitted $H$ (10--1000 s)} & \multicolumn{3}{c}{$\Delta H$} & \multicolumn{3}{c}{$\Delta H_\mathrm{loc}(1000~\mathrm{s})$} \\
\cmidrule(lr){3-5}\cmidrule(lr){6-8}\cmidrule(l){9-11}
compared within & \makecell{G\\kept (\%)} & $\G$ & $\Sl\mid\G$ & $\A\mid\G$ & \makecell{$\G-\Sl$\\as is} & \makecell{$\G-\Sl$\\reweighted} & \makecell{$\G-\A$\\reweighted} & \makecell{$\G-\Sl$\\as is} & \makecell{$\G-\Sl$\\reweighted} & \makecell{$\G-\A$\\reweighted} \\
\midrule
none & 100 & 0.916 & 0.899 & 0.914 & $+$0.017 & \makecell{$+$0.017\\{\scriptsize[$+$0.016, $+$0.018]}} & \makecell{$+$0.002\\{\scriptsize[$+$0.001, $+$0.003]}} & $+$0.055 & \makecell{$+$0.055\\{\scriptsize[$+$0.053, $+$0.057]}} & \makecell{$+$0.032\\{\scriptsize[$+$0.030, $+$0.034]}} \\[2pt]
zone & 100 & 0.919 & 0.902 & 0.916 & $+$0.021 & \makecell{$+$0.018\\{\scriptsize[$+$0.017, $+$0.019]}} & \makecell{$+$0.003\\{\scriptsize[$+$0.003, $+$0.004]}} & $+$0.059 & \makecell{$+$0.056\\{\scriptsize[$+$0.054, $+$0.058]}} & \makecell{$+$0.028\\{\scriptsize[$+$0.026, $+$0.029]}} \\[2pt]
zone $\times$ month & 99 & 0.919 & 0.902 & 0.915 & $+$0.022 & \makecell{$+$0.017\\{\scriptsize[$+$0.016, $+$0.018]}} & \makecell{$+$0.003\\{\scriptsize[$+$0.003, $+$0.004]}} & $+$0.064 & \makecell{$+$0.054\\{\scriptsize[$+$0.053, $+$0.057]}} & \makecell{$+$0.028\\{\scriptsize[$+$0.027, $+$0.031]}} \\[2pt]
zone $\times$ hour & 99 & 0.919 & 0.902 & 0.913 & $+$0.022 & \makecell{$+$0.017\\{\scriptsize[$+$0.016, $+$0.018]}} & \makecell{$+$0.005\\{\scriptsize[$+$0.005, $+$0.006]}} & $+$0.064 & \makecell{$+$0.052\\{\scriptsize[$+$0.050, $+$0.054]}} & \makecell{$+$0.026\\{\scriptsize[$+$0.025, $+$0.028]}} \\[2pt]
site-day $\times$ hour & 53 & 0.917 & 0.904 & 0.913 & $+$0.015 & \makecell{$+$0.013\\{\scriptsize[$+$0.012, $+$0.014]}} & \makecell{$+$0.008\\{\scriptsize[$+$0.007, $+$0.009]}} & $+$0.051 & \makecell{$+$0.044\\{\scriptsize[$+$0.041, $+$0.046]}} & \makecell{$+$0.033\\{\scriptsize[$+$0.031, $+$0.036]}} \\[2pt]
\bottomrule
\end{tabular}

\caption{Fitted $H$ and differences by stratification. ``Reweighted'': the other population reweighted to the strata of G (value, 5--95\% band);
``as is'': the same strata without reweighting. $H_{\Sl\mid\G}$: $H$ of S reweighted to G.}
\label{tab:zones}
\end{table}

\begin{reading}
\textbf{Reading.}
\begin{itemize}
  \item \textbf{The gap opens above $\tau\approx150$~s.} $\Delta H_\mathrm{loc}(\tau)$ between G and S is small below it ($\ZoneShortMin$ to $\ZoneShortMax$ from 20 to 140~s)
        and rises to \ZoneDHlocSNone{} at 1000~s, above one climb--glide cycle: the group acts on how cycles are chained, not on a single glide or climb.
  \item \textbf{The zone does not explain it.} Reweighting S to the cells where G flies lowers $\Delta H_\mathrm{loc}(1000)$ by only \ZoneRedLocZone\%
        (\ZoneDHlocSupportZone{} $\to$ \ZoneDHlocSZone) and the fitted $\Delta H$ by \ZoneRedHZone\%; with the month, \ZoneRedLocZoneMonth\% and \ZoneRedHZoneMonth\%; with the
        hour, \ZoneRedLocZoneHour\% and \ZoneRedHZoneHour\%. At least 80\% of the long-lag gap remains. Solo pilots fly where groups fly: \ZoneKeptZone\% of G has a solo
        counterpart in its cell.
  \item \textbf{Same day, site and hour} (the finest strata, \ZoneKeptSdh\% of G kept): $\Delta H_\mathrm{loc}(1000)=\ZoneDHlocSSdh$ \ZoneDHlocSSdhBand, \ZoneRedLocSdh\% below its as-is value;
        fitted $\Delta H=\ZoneDHSSdh$ \ZoneDHSSdhBand. Solo flights on the same day, site and hour as a group still progress less straight over 1000~s.
  \item \textbf{Fitted $H$:} G \ZoneHNoneG{}, S \ZoneHNoneS{} raw. S reweighted to the zones of G: \ZoneHZoneS{}; to zone$\times$hour: \ZoneHZoneHourS{}; to site-day$\times$hour: \ZoneHSdhS{} (G \ZoneHSdhG).
        Reweighting raises the fitted $H$ of S by a few thousandths and never brings it to that of G.
  \item \textbf{Against the same flights alone} (A): $\Delta H_\mathrm{loc}(1000)$ is \ZoneDHlocANone{} pooled and \ZoneDHlocAZone, \ZoneDHlocAZoneHour, \ZoneDHlocASdh{} in zone, zone$\times$hour,
        site-day$\times$hour strata: the same pilots in the same zones and hours progress straighter in a group. Below 150~s the sign depends on the hour: A is above G pooled and equal at the same site-day and hour,
        which is why the within-flight estimate of Section~\ref{sec:other} must be read with the same-site-day-hour one.
  \item \textbf{Answer:} the difference between group and solo flights does not come from the zones where groups fly, nor from the month or hour at that place;
        it is reduced by about a fifth by the day, site and hour together, and the rest stays with the pilot and the day fixed. It is not a zone effect.
\end{itemize}
\textbf{Limits.} Cells are 25~km, from the middle of a piece; the terrain inside a cell (ridge, plain) is not controlled, and a group may choose better air within it:
this is part of what we call the group effect. Paragliders only.
\end{reading}
